\pagestyle{empty}
\tight
\begin{tblr}{width=\textwidth,colspec={|Q[c,m,1.9cm]|X[l,m]|},hlines,vlines,hspan=minimal,row{1}={bg=headblue},rowsep=1pt,colsep=3pt}
\SetCell{c,m}\hfil\logo\hfil & \textbf{POLITEKNIK NEGERI MALANG}\par\textbf{JURUSAN TEKNOLOGI INFORMASI}\par\textbf{PROGRAM STUDI: D4 TEKNIK INFORMATIKA} \\
\SetCell[c=2]{c} \textbf{RENCANA PEMBELAJARAN SEMESTER (RPS)} & \\
\end{tblr}
\begin{longtblr}{width=\textwidth,colspec={|Q[l,m,1.9cm]|X[0.85,l,m]|X[1.45,l,m]|X[0.75,l,m]|X[1.15,l,m]|},hlines,vlines,hspan=minimal,rowsep=1pt,colsep=2pt,row{1,3}={bg=headgray,font=\bfseries}}
MATA KULIAH & KODE & BOBOT (SKS)/jam & SEMESTER & TGL. PENYUSUNAN \\
Matematika Dasar & RTI251004 & 2 SKS/4 jam & 1 & 7 Agustus 2025 \\
OTORISASI & Dosen Pengembang RPS & Koordinator RMK & \SetCell[c=2]{l} Ka PRODI &  \\
 & Prof. Dr. Eng. Cahya Rahmad, S.T., M.Kom. &  & \SetCell[c=2]{l}{\vspace{8mm}\par Dr. Ely Setyo Astuti, S.T., M.T.} &  \\
\SetCell[r=9]{l} Capaian Pembelajaran (CP) & \SetCell[c=4]{l,bg=headgray,font=\bfseries} Capaian Pembelajaran Lulusan Program Studi (CPL-Prodi) &  &  &  \\
 & CPL09 & \SetCell[c=3]{l} Mampu menjelaskan cara kerja sistem komputer dan menerapkan pendekatan komputasi untuk menyelesaikan masalah. &  &  \\
 & \SetCell[c=4]{l,bg=headgray,font=\bfseries} Capaian Pembelajaran mata kuliah (CPL-MK) &  &  &  \\
 & CPMK0902 & \SetCell[c=3]{l} Mahasiswa mampu menerapkan prinsip matematis dan struktur diskrit untuk analisis komputasional. &  &  \\
 & \SetCell[c=4]{l,bg=headgray,font=\bfseries} Kemampuan akhir tiap tahapan belajar (Sub-CPMK) &  &  &  \\
 & SCPMK0902-00401 & \SetCell[c=3]{l} Mahasiswa mampu menerapkan logika proposisi dan operasi himpunan untuk memodelkan relasi dan menyusun solusi masalah komputasional sederhana. &  &  \\
 & SCPMK0902-00402 & \SetCell[c=3]{l} Mahasiswa mampu melakukan perhitungan pada relasi, fungsi, dan konversi antarsistem bilangan (biner, oktal, desimal, heksadesimal). &  &  \\
 & SCPMK0902-00403 & \SetCell[c=3]{l} Mahasiswa mampu menerapkan induksi matematika, aljabar Boolean, dan kombinatorial dalam analisis persoalan komputasional. &  &  \\
 & SCPMK0902-00404 & \SetCell[c=3]{l} Mahasiswa mampu menerapkan konsep graf dan pohon untuk merepresentasikan dan menyelesaikan persoalan komputasi. &  &  \\
\SetCell[r=2]{l} Penilaian & \SetCell[c=4]{l} *Nilai CPMK didapat dari agregasi nilai SUB-CPMK yang dibebankan &  &  &  \\
 & \SetCell[c=4]{l}{\tiny\begin{tblr}{width=\linewidth,colspec={|Q[c,m,0.1\linewidth]|Q[c,m,0.16\linewidth]|X[l,m]|X[l,m]|X[l,m]|X[l,m]|X[l,m]|Q[c,m,0.08\linewidth]|},hlines,vlines,hspan=minimal,rowsep=1pt,colsep=0.7pt,row{1,2}={bg=headgray,font=\bfseries}}
Kode CPMK & Kode SCPMK & \SetCell[c=5]{c} Bobot per Bentuk Penilaian & & & & & Total Bobot \\
 & & Tugas & Kuis 1 & UTS & Kuis 2 & UAS & \\
\SetCell[r=1]{c} \SetCell[r=4]{c} CPMK0902 & SCPMK0902-00401 & 5\%: latihan soal logika dan himpunan & 7.5\%: logika dan himpunan & 15\%: logika dan himpunan & & & 27.5\% \\
\SetCell[r=3]{c}  & SCPMK0902-00402 & 4.5\%: latihan soal relasi, fungsi, sistem bilangan & & 15\%: relasi, fungsi, dan sistem bilangan & & & 19.5\% \\
 & SCPMK0902-00403 & 6\%: latihan soal induksi, Boolean, kombinatorial & & & 7.5\%: induksi, Boolean, dan kombinatorial & 17.5\%: induksi, Boolean, dan kombinatorial & 31\% \\
 & SCPMK0902-00404 & 4.5\%: latihan soal graf dan pohon & & & & 17.5\%: graf dan pohon & 22\% \\
\SetCell[c=7]{r}\textbf{Total Per Penilaian} & & & & & & & \textbf{100\%} \\
\end{tblr}} &  &  &  \\
Diskripsi Singkat Mata Kuliah & \SetCell[c=4]{l} Mata kuliah Matematika Dasar mempelajari matematika diskrit, yaitu cabang aljabar yang mengkaji perhitungan bilangan diskrit (tidak kontinyu) dan cara paling efisien menemukan solusinya; dalam konteks teknologi informasi, bilangan biner mendapatkan perhatian utama. &  &  &  \\
Bahan Kajian / Materi Pembelajaran & \SetCell[c=4]{l}\begin{enumerate}\item Proposisi dan logika\item Teori himpunan\item Relasi dan fungsi\item Teori bilangan dan sistem bilangan\item Induksi dan rekursi\item Aljabar Boolean\item Kombinatorial\item Teori graf dan pohon\end{enumerate} &  &  &  \\
\SetCell[r=2]{l} Pustaka & \SetCell[c=4]{l} \textbf{Utama:}\begin{enumerate}\item Yan Watequlis, Cahya Rahmad, Deasy Sandhya Elya. 2017. Matematika Diskrit. Polinema Press.\end{enumerate} &  &  &  \\
 & \SetCell[c=4]{l} \textbf{Pendukung:}\begin{enumerate}\item Steven G. Krantz. 2009. Discrete Mathematics Demystified. McGraw-Hill.\item Kenneth H. Rosen. 1999. Discrete Mathematics and Its Application. McGraw-Hill.\item C.L. Liu. 1985. Element of Discrete Mathematics. McGraw-Hill, Inc.\item Munir, Rinaldi. 2012. Matematika Diskrit Ed. Revisi Ke-5. Informatika Bandung.\end{enumerate} &  &  &  \\
Media Pembelajaran & \SetCell[c=2]{l} \textbf{Software:} LMS dan office suite. &  & \SetCell[c=2]{l} \textbf{Hardware:} Komputer/laptop, proyektor. &  \\
Nama Dosen Pengampu & \SetCell[c=4]{l}\begin{enumerate}\item Muhammad Afif Hendrawan, S.Kom., M.T.\item Prof. Dr. Eng. Cahya Rahmad, S.T., M.Kom.\item Habibie Ed Dien, S.Kom., M.T.\item Erfan Rohadi, S.T., M.Eng., Ph.D.\item Vipkas Al Hadid Firdaus, S.T., M.T.\end{enumerate} &  &  &  \\
Matakuliah Syarat & \SetCell[c=4]{l} - &  &  &  \\
\end{longtblr}
\begin{longtblr}{width=\textwidth,colspec={|Q[c,m,0.06\textwidth]|X[1.35,l,m]|X[1.08,l,m]|X[1.16,l,m]|X[0.7,l,m]|X[1.24,l,m]|X[1.06,l,m]|X[0.98,l,m]|Q[c,m,0.05\textwidth]|},hlines,vlines,hspan=minimal,rowhead=2,row{1,2}={bg=headgreen,font=\bfseries},rowsep=1pt,colsep=1.5pt}
Mgg.\par Ke & Kemampuan Akhir Yang Direncanakan (Sub-CP-MK) & Bahan kajian (Materi Pembelajaran) & Bentuk dan Metode Pembelajaran & Estimasi Waktu & Pengalaman Belajar Mahasiswa & Kriteria \& Bentuk Penilaian & Indikator Penilaian & Bobot Penilaian (\%) \\
(1) & (2) & (3) & (4) & (5) & (6) & (7) & (8) & (9) \\
1 & SCPMK0902-00401 & Konsep diskrit; proposisi; proposisi majemuk; ekivalensi, tautologi, dan kontradiksi. & Bentuk: Kuliah luring/daring. Metode: Diskusi kelas dan kelompok, demo contoh soal. Penugasan: latihan soal logika. & 1 x 2 x 50' tatap muka; 1 x 2 x 50' latihan/tugas mandiri. & Mahasiswa memahami pengantar matematika diskrit dan menentukan nilai kebenaran proposisi melalui latihan soal. & Bentuk penilaian: Tugas latihan soal. Kriteria: ketepatan penjelasan dan jawaban. Mengacu rubrik RTM. & Ketepatan menjelaskan konsep diskrit dan menentukan nilai kebenaran proposisi majemuk. & 2 \\
2 & SCPMK0902-00401 & Hukum-hukum logika; hukum De Morgan; proposisi bersyarat; proposisi bikondisional (dwisyarat). & Bentuk: Kuliah luring/daring. Metode: Diskusi kelas dan kelompok, demo contoh soal. Penugasan: latihan soal logika. & 1 x 2 x 50' tatap muka; 1 x 2 x 50' latihan/tugas mandiri. & Mahasiswa menyelesaikan latihan soal hukum logika dan mengerjakan soal secara spontan di papan tulis. & Bentuk penilaian: Tugas latihan soal. Kriteria: ketepatan penjelasan dan jawaban. Mengacu rubrik RTM. & Ketepatan menerapkan hukum logika dan De Morgan dalam penyederhanaan proposisi. & 1.5 \\
3 & SCPMK0902-00401 & Definisi himpunan; himpunan kosong; kardinalitas; subset; operasi himpunan; hukum-hukum himpunan; prinsip inklusi-eksklusi. & Bentuk: Kuliah luring/daring. Metode: Diskusi kelas dan kelompok, demo contoh soal. Penugasan: latihan soal himpunan. & 1 x 2 x 50' tatap muka; 1 x 2 x 50' latihan/tugas mandiri. & Mahasiswa menyelesaikan latihan soal operasi himpunan dan diagram Venn serta membahasnya di kelas. & Bentuk penilaian: Tugas latihan soal. Kriteria: ketepatan penjelasan dan jawaban. Mengacu rubrik RTM. & Ketepatan melakukan operasi himpunan dan menerapkan prinsip inklusi-eksklusi. & 1.5 \\
4 & SCPMK0902-00401 & Kuis 1: materi pertemuan 1 s.d. 3 (logika dan himpunan). & Bentuk: Ujian tulis daring/luring. Metode: mengerjakan soal secara individual. & 1 x 2 x 50' ujian; 1 x 2 x 50' reviu mandiri. & Mahasiswa menjawab soal kuis logika dan himpunan dengan tepat dan mandiri. & Bentuk penilaian: Kuis 1. Kriteria: ketepatan jawaban. Mengacu rubrik RTM. & Ketepatan jawaban soal logika proposisi dan operasi himpunan. & 7.5 \\
5 & SCPMK0902-00402 & Definisi relasi; jenis-jenis relasi; definisi fungsi; fungsi eksponensial dan logaritmis; fungsi rekursif; kardinalitas. & Bentuk: Kuliah luring/daring. Metode: Diskusi kelas dan kelompok, demo contoh soal. Penugasan: latihan soal relasi dan fungsi. & 1 x 2 x 50' tatap muka; 1 x 2 x 50' latihan/tugas mandiri. & Mahasiswa mengerjakan latihan soal relasi dan fungsi serta soal take-home penerapannya dalam kehidupan sehari-hari. & Bentuk penilaian: Tugas latihan soal. Kriteria: ketepatan penjelasan dan jawaban. Mengacu rubrik RTM. & Ketepatan menentukan jenis relasi dan menghitung nilai fungsi. & 1.5 \\
6 & SCPMK0902-00402 & Penulisan baku sistem bilangan; konversi bilangan biner dan oktal. & Bentuk: Kuliah luring/daring. Metode: Diskusi kelas dan kelompok, demo contoh soal. Penugasan: latihan soal sistem bilangan. & 1 x 2 x 50' tatap muka; 1 x 2 x 50' latihan/tugas mandiri. & Mahasiswa menyelesaikan latihan soal konversi bilangan biner dan oktal serta membahasnya di papan tulis. & Bentuk penilaian: Tugas latihan soal. Kriteria: ketepatan penjelasan dan jawaban. Mengacu rubrik RTM. & Ketepatan melakukan konversi bilangan biner dan oktal. & 1.5 \\
7 & SCPMK0902-00402 & Konversi bilangan desimal dan heksadesimal; operasi bilangan. & Bentuk: Kuliah luring/daring. Metode: Diskusi kelas dan kelompok, demo contoh soal. Penugasan: latihan soal sistem bilangan. & 1 x 2 x 50' tatap muka; 1 x 2 x 50' latihan/tugas mandiri. & Mahasiswa menyelesaikan latihan soal konversi desimal-heksadesimal dan operasi antarbasis bilangan. & Bentuk penilaian: Tugas latihan soal. Kriteria: ketepatan penjelasan dan jawaban. Mengacu rubrik RTM. & Ketepatan konversi desimal-heksadesimal dan operasi antarbasis bilangan. & 1.5 \\
8 & SCPMK0902-00401, SCPMK0902-00402 & UTS: materi pertemuan 1 s.d. 7. & Bentuk: Ujian tulis daring/luring. Metode: mengerjakan soal secara individual. & 1 x 2 x 50' ujian; 1 x 2 x 50' reviu mandiri. & Mahasiswa menjawab soal UTS logika, himpunan, relasi, fungsi, dan sistem bilangan dengan tepat. & Bentuk penilaian: UTS. Kriteria: ketepatan jawaban. Mengacu rubrik RTM. & Ketepatan jawaban soal logika, himpunan, relasi, fungsi, dan sistem bilangan. & 30 \\
9 & SCPMK0902-00403 & Prinsip induksi matematika; pembuktian dengan induksi matematika. & Bentuk: Kuliah luring/daring. Metode: Diskusi kelas dan kelompok, demo contoh soal. Penugasan: latihan soal induksi. & 1 x 2 x 50' tatap muka; 1 x 2 x 50' latihan/tugas mandiri. & Mahasiswa menyusun pembuktian induksi matematika melalui latihan soal terbimbing dan mandiri. & Bentuk penilaian: Tugas latihan soal. Kriteria: ketepatan penjelasan dan jawaban. Mengacu rubrik RTM. & Ketepatan menyusun langkah basis dan langkah induksi pada pembuktian. & 1.5 \\
10 & SCPMK0902-00403 & Teorema dasar aljabar Boolean; membuat fungsi dari tabel; menyederhanakan fungsi Boolean. & Bentuk: Kuliah luring/daring. Metode: Diskusi kelas dan kelompok, demo contoh soal. Penugasan: latihan soal aljabar Boolean. & 1 x 2 x 50' tatap muka; 1 x 2 x 50' latihan/tugas mandiri. & Mahasiswa menyusun dan menyederhanakan fungsi Boolean dari tabel kebenaran melalui latihan soal. & Bentuk penilaian: Tugas latihan soal. Kriteria: ketepatan penjelasan dan jawaban. Mengacu rubrik RTM. & Ketepatan menyusun dan menyederhanakan fungsi Boolean. & 1.5 \\
11 & SCPMK0902-00403 & Kaidah penjumlahan dan perkalian; kombinatorial dasar. & Bentuk: Kuliah luring/daring. Metode: Diskusi kelas dan kelompok, demo contoh soal. Penugasan: latihan soal kombinatorial. & 1 x 2 x 50' tatap muka; 1 x 2 x 50' latihan/tugas mandiri. & Mahasiswa mengerjakan latihan soal kombinatorial dasar dan soal take-home penerapannya sehari-hari. & Bentuk penilaian: Tugas latihan soal. Kriteria: ketepatan penjelasan dan jawaban. Mengacu rubrik RTM. & Ketepatan menerapkan kaidah penjumlahan dan perkalian pada pencacahan. & 1.5 \\
12 & SCPMK0902-00403 & Permutasi; kombinasi; kombinasi dengan perulangan objek. & Bentuk: Kuliah luring/daring. Metode: Diskusi kelas dan kelompok, demo contoh soal. Penugasan: latihan soal kombinatorial. & 1 x 2 x 50' tatap muka; 1 x 2 x 50' latihan/tugas mandiri. & Mahasiswa mengerjakan latihan soal permutasi dan kombinasi serta membahas solusinya di kelas. & Bentuk penilaian: Tugas latihan soal. Kriteria: ketepatan penjelasan dan jawaban. Mengacu rubrik RTM. & Ketepatan memilih dan menghitung rumus permutasi dan kombinasi. & 1.5 \\
13 & SCPMK0902-00403 & Kuis 2: materi pertemuan 9 s.d. 12 (induksi, aljabar Boolean, kombinatorial). & Bentuk: Ujian tulis daring/luring. Metode: mengerjakan soal secara individual. & 1 x 2 x 50' ujian; 1 x 2 x 50' reviu mandiri. & Mahasiswa menjawab soal kuis induksi, aljabar Boolean, dan kombinatorial dengan tepat dan mandiri. & Bentuk penilaian: Kuis 2. Kriteria: ketepatan jawaban. Mengacu rubrik RTM. & Ketepatan jawaban soal induksi, aljabar Boolean, dan kombinatorial. & 7.5 \\
14 & SCPMK0902-00404 & Definisi graf; jenis graf; terminologi graf. & Bentuk: Kuliah luring/daring. Metode: Diskusi kelas dan kelompok, demo contoh soal. Penugasan: latihan soal graf. & 1 x 2 x 50' tatap muka; 1 x 2 x 50' latihan/tugas mandiri. & Mahasiswa mengerjakan latihan soal representasi dan terminologi graf serta membahasnya di kelas. & Bentuk penilaian: Tugas latihan soal. Kriteria: ketepatan penjelasan dan jawaban. Mengacu rubrik RTM. & Ketepatan mengidentifikasi jenis graf dan menggunakan terminologi graf. & 1.5 \\
15 & SCPMK0902-00404 & Lintasan dan sirkuit Euler; lintasan dan sirkuit Hamilton. & Bentuk: Kuliah luring/daring. Metode: Diskusi kelas dan kelompok, demo contoh soal. Penugasan: latihan soal graf. & 1 x 2 x 50' tatap muka; 1 x 2 x 50' latihan/tugas mandiri. & Mahasiswa menentukan lintasan dan sirkuit Euler/Hamilton pada graf melalui latihan soal. & Bentuk penilaian: Tugas latihan soal. Kriteria: ketepatan penjelasan dan jawaban. Mengacu rubrik RTM. & Ketepatan menentukan lintasan dan sirkuit Euler serta Hamilton. & 1.5 \\
16 & SCPMK0902-00404 & Definisi pohon; spanning tree; pohon berakar (rooted tree); pohon terurut (ordered tree); pohon n-ary; pohon biner (binary tree). & Bentuk: Kuliah luring/daring. Metode: Diskusi kelas dan kelompok, demo contoh soal. Penugasan: latihan soal pohon. & 1 x 2 x 50' tatap muka; 1 x 2 x 50' latihan/tugas mandiri. & Mahasiswa mengerjakan latihan soal pohon dan soal take-home penerapan pohon dalam kehidupan sehari-hari. & Bentuk penilaian: Tugas latihan soal. Kriteria: ketepatan penjelasan dan jawaban. Mengacu rubrik RTM. & Ketepatan membangun spanning tree dan menganalisis pohon biner. & 1.5 \\
17 & SCPMK0902-00403, SCPMK0902-00404 & UAS: materi pertemuan 1 s.d. 16 dengan penekanan induksi, aljabar Boolean, kombinatorial, graf, dan pohon. & Bentuk: Ujian tulis daring/luring. Metode: mengerjakan soal secara individual. & 1 x 2 x 50' ujian; 1 x 2 x 50' reviu mandiri. & Mahasiswa mengerjakan soal UAS dengan baik dan menunjukkan langkah penyelesaian yang runtut. & Bentuk penilaian: UAS. Kriteria: ketepatan jawaban. Mengacu rubrik RTM. & Ketepatan jawaban soal induksi, aljabar Boolean, kombinatorial, graf, dan pohon. & 35 \\
\end{longtblr}
